\begin{resumo}[Abstract]
\begin{otherlanguage*}{english}
This work is situated in the context of the continuous growth of digital data generation and the search for long-term storage alternatives with higher density and lower energy demand than conventional media. The objective of this study was to develop and evaluate a modular simulator of the DNA Fountain pipeline, capable of representing the encoding, biological screening, erasure channel, and peeling decoding stages, with emphasis on LT (\emph{Luby Transform}) codes, the robust soliton distribution (RSD), and the \emph{peeling} decoding algorithm. The methodology consists of formalizing the LT encoding model, representing the decoding process through bipartite graphs over the binary finite field \(GF(2)\), and implementing a pure Python modular reference simulator. The simulator generates controlled DNA Fountain instances, applies biological screening constraints, models probabilistic erasure channels, and executes the decoding process to evaluate data reconstruction over a factorial grid of 216 parameter combinations with 100 repetitions ($21,600$ trials). Results reveal that, for finite block regimes ($K \le 64$), appropriate parameterizations ($c = 0.1$, $\delta = 0.10$) achieved 100\% decoding success for $K = 32$ and $K = 64$ with overhead $\epsilon = 1.00$ over ideal channels (95\% exact Clopper--Pearson confidence interval between $96.4\%$ and $100.0\%$, and 90\% for $K = 16$ due to ripple die-out from lack of degree-1 equations). Furthermore, for the configurations evaluated with $K = 32$, $c \ge 0.5$ concentrated probability mass on unit degrees and resulted in $0\%$ decoding success ($\epsilon = 0.50$ and $\epsilon = 1.00$ over ideal channel), while sustaining $82\%$ success for $K = 64$ due to the shift of the $\tau(d)$ peak to degree 2. Additionally, biological screening rejected approximately 40\% of candidate oligos with limited differences in sampled mean degree, sustaining resilience under up to 20\% erasure channel losses given sufficient redundancy. This work is delimited as a Software Engineering investigation applied to the modeling, simulation, and experimental analysis of fountain codes for DNA data storage, without including chemical synthesis, sequencing, or physical laboratory validation.

\vspace{\onelineskip}

\noindent
\textbf{Keywords}: DNA data storage; DNA Fountain; LT codes; \emph{peeling} decoding; software simulation.
\end{otherlanguage*}
\end{resumo}
